\begin{theorem}[Conjecture: no revivals outside the admissible set]
\label{conj:no-other-k}
No cycle length $k \notin \{2, 3, 4, 5, 6, 8, 10\}$ admits a
non-trivial 1D full-state revival at any period $N$ or any coin
parameter $\rho \in (0, 1)$.
\end{theorem}

\begin{proof}[Sketch]
The dispersion relation for a $k$-cycle is
$\sin\theta_j = \sqrt{\rho}\,\sin(2\pi j / k)$ for
$j = 1, \ldots, k-1$. A revival at period $N$ forces each
eigenphase to satisfy $\theta_j = \pi r_j / N$ for some integer
$r_j$. Taking the ratio of the dispersion relations for the two
lowest spatial modes ($j = 2$ and $j = 1$) eliminates $\rho$
entirely:
\begin{equation}
\frac{\sin\theta_2}{\sin\theta_1}
= \frac{\sin(4\pi / k)}{\sin(2\pi / k)}
= 2\cos\!\left(\frac{2\pi}{k}\right).
\label{eq:mode-ratio}
\end{equation}
Because $\theta_1, \theta_2 \in \pi\mathbb{Q}$, the left-hand
side of Eq.~\eqref{eq:mode-ratio} lies in the real subfield
$\mathbb{Q}(\zeta_{2N})^+$ of the cyclotomic field generated by
$2N$-th roots of unity. The right-hand side lies in the real
subfield
$\mathbb{Q}(\zeta_k)^+ = \mathbb{Q}(\cos(2\pi / k))$, whose
degree over $\mathbb{Q}$ is $\phi(k)/2$ for $k \ge 3$ (see,
e.g., Ref.~\cite{Dukes2014} for the relevant algebraic
background). This produces a necessary condition: the field
$\mathbb{Q}(\cos(2\pi/k))$ must embed in
$\mathbb{Q}(\zeta_{2N})^+$, which restricts $N$ to a sparse
arithmetic set determined by the intersection of the two
cyclotomic fields. Two cases are tractable. For $k = 12$ the
mode ratio forces $\sqrt{\rho} = \sin(\pi/2) = 1$, collapsing
the coin to a deterministic reflector. For $k = 7$ and $k = 9$
the field $\mathbb{Q}(\cos(2\pi/k))$ is cubic, and the
embedding condition restricts $N$ to multiples of $k$; a finite
check over these $N$ would close the argument in each case.
The general case, $k \ge 11$ with $\phi(k)/2 \ge 3$, requires a
systematic treatment of the cyclotomic intersection
$\mathbb{Q}(\cos(2\pi/k)) \cap \mathbb{Q}(\zeta_{2N})^+$ that is
not currently available. We therefore state the conclusion as a
conjecture, supported on the search bound of
Theorem~\ref{thm:1d-nonexistence} and consistent with the
explicit witnesses of Theorem~\ref{thm:1d-existence}.
\end{proof}
